% ============================================================
%  INTRODUCTION SECTION
%  Quark Structure of Hadrons / Particle Physics Project
% ============================================================

\section{Introduction}
\label{sec:introduction}

The matter that constitutes the observable universe is composed, at the most
fundamental level, of two families of spin-$\tfrac{1}{2}$ fermions: quarks and
leptons. Together with the gauge bosons that mediate the three non-gravitational
forces, these elementary constituents form the content of the Standard Model of
particle physics. Leptons interact through the electromagnetic and
weak forces but are oblivious to the strong interaction. Quarks, by contrast,
carry colour charge and therefore participate in all three interactions; it is
precisely the strong force that binds them into the composite, strongly interacting
particles collectively known as \emph{hadrons} \cite[Ch.~1]{HalzenMartin1984}.

%% FIGURE SUGGESTION: Standard Model particle table (leptons, quarks, gauge bosons)
%% e.g.\ the canonical SM table from PDG or from Halzen \& Martin Ch.~1.

Hadrons are classified according to their spin into two broad categories. Particles
with half-integer total angular momentum are called \emph{baryons}, while those with integer spin are called \emph{mesons}. The distinction carries a deeper dynamical
significance: baryon number $B$ is an additive, conserved quantum number, so
that baryons and antibaryons possess $B = +1$ and $B = -1$ respectively, and
baryon number is conserved in every known interaction. Mesons carry $B = 0$ and
can therefore be freely created or annihilated, subject only to energy conservation
and charge conservation \cite[Ch.~1]{Amsler2018}.

\subsection{Historical Background and the Quark Model}

Prior to the 1960s, advances in accelerator technology unleashed an avalanche of
newly discovered strongly interacting particles---protons, neutrons, pions, kaons,
hyperons, and their numerous excited states---that earned the era the description
``the particle zoo''. Organizing this proliferation demanded a classification
principle analogous to Mendeleev's periodic table. A pivotal step came with the
recognition that hadrons cluster naturally into multiplets of the approximate
flavour symmetry group $\mathrm{SU}(3)_f$, proposed independently around
1961--1964 by Gell-Mann and Ne'eman under the name \emph{the Eightfold Way}
\cite[Ch.~2]{HalzenMartin1984}. This symmetry groups the lightest pseudoscalar and
vector mesons into octets (and singlets that mix to form nonets), and arranges the
lightest baryons into an octet and a decuplet, a scheme whose predictive success
was spectacularly confirmed by the observation of the $\Omega^-$ baryon in 1964 at
Brookhaven, with a mass strikingly close to the prediction derived from the pattern
of the baryon decuplet \cite[Ch.~7]{Amsler2018}.

The natural explanation of these multiplet structures emerged in 1964 when
Gell-Mann and, independently, Zweig proposed that hadrons are not elementary but
are built from more fundamental constituents. Gell-Mann named these constituents
\emph{quarks}, drawing on a phrase in James Joyce's \textit{Finnegans Wake}:
``Three quarks for Muster Mark''. In Gell-Mann's own cautious words at the time,
``such particles presumably are not real, but we may use them in our field theory
anyway'' \cite[Ch.~1]{Amsler2018}. The subsequent decades proved this modesty
unnecessary: deep inelastic scattering experiments at SLAC in the late 1960s
revealed point-like constituents inside the proton---the \emph{partons} of
Feynman's 1969 description---which were eventually identified with the quarks of
the model, together with electrically neutral gluons that carry a share of the
proton's momentum. The observation of narrow two- and three-jet events at PETRA
and PEP provided direct visual evidence for quark and gluon fragmentation into
hadronic matter \cite[Ch.~14]{HalzenMartin1984}.

\subsection{QCD and the strong interaction}

\subsection{Screening and anti-screening}

\subsection{Reminder: scattering}

\subsection{Yukawa potential and massive carriers: the remanent nuclear forces}

\subsection{Parity}

